% ---------------------------------------------------------------------------
% Section 2: the model and its exact free-fermion duality
% ---------------------------------------------------------------------------
% Section-local macros (unique names, \providecommand only).
\providecommand{\sxb}[1]{\sigma^{x}_{b_{#1}}}
\providecommand{\szb}[1]{\sigma^{z}_{b_{#1}}}
\providecommand{\taux}[1]{\tau^{x}_{#1}}
\providecommand{\tauz}[1]{\tau^{z}_{#1}}
\providecommand{\cdr}{\tilde c}
\providecommand{\LinkS}{\mathcal{S}}
\providecommand{\SiteT}{\mathcal{T}}

\section{The model and its exact free-fermion duality}\label{sec:model}

In this section we define the model, prove the exact duality to free fermions coupled to static
$\Z_2$ charges (\cref{thm:duality}), show that the duality is local across a cut up to a single
controlled gate (\cref{lem:cutlocal}), and derive the exact sector-averaging formulas for the
initial states of \cite{SKKM2017} (\cref{cor:sector}) and for general product initial states
(\cref{lem:productlaw}).  Apart from \cref{rem:gaussproj}(iii), which concerns $\ell^2(\Z)$, everything
in this section is an exact identity in finite volume; no probabilistic input and no localization input is
used.  The one-body input enters only in \cref{sec:onebody}.  The duality itself is the one of
\cite{SKKM2017} (\cref{rem:boundary}(b)); what we add is a convention (the boundary link) under which it
is an exact unitary equivalence on the full Hilbert space, without constraints, together with the
precise bookkeeping of signs, boundary variants, cuts and initial states that the later sections need.

\subsection{Hilbert space, Hamiltonian and Gauss law}\label{sec:model:def}

\paragraph{Matter.}
Fix $N\ge 1$ and the sites $j\in\{1,\dots,N\}$.  Let $\cF_N$ be the fermionic Fock space over
$\C^N=\ell^2(\{1,\dots,N\})$, with annihilation operators $f_1,\dots,f_N$,
$\{f_i,f_j^\dagger\}=\delta_{ij}$, $\{f_i,f_j\}=0$, and $n_j=f_j^\dagger f_j$.  For
$X\subset\{1,\dots,N\}$ we write $N_X=\sum_{j\in X}n_j$, $N_f=N_{\{1,\dots,N\}}$,
$N_{\le j}=N_{\{1,\dots,j\}}$, $N_{<j}=N_{\{1,\dots,j-1\}}$, $N_{>j}=N_{\{j+1,\dots,N\}}$.
The occupation basis is fixed once and for all by the ordering convention
\begin{equation}\label{eq:occbasis}
  \ket{n}:=(f_1^\dagger)^{n_1}(f_2^\dagger)^{n_2}\cdots(f_N^\dagger)^{n_N}\ket{0},\qquad
  n=(n_1,\dots,n_N)\in\{0,1\}^N ,
\end{equation}
so that, with $e_j$ the $j$-th unit vector and $N_{<j}(n)=\sum_{k<j}n_k$,
\begin{equation}\label{eq:fsign}
  f_j\ket{n}=\delta_{n_j,1}\,(-1)^{N_{<j}(n)}\ket{n-e_j},\qquad
  f_j^\dagger\ket{n}=\delta_{n_j,0}\,(-1)^{N_{<j}(n)}\ket{n+e_j}.
\end{equation}
This is the Jordan--Wigner ordering ``site $1$ first''.  An operator is \emph{even} if it commutes
with the fermion parity $(-1)^{N_f}$ and \emph{odd} if it anticommutes with it.

\paragraph{Gauge field.}
The links are $b_j=(j,j+1)$, $j=1,\dots,N$.  Each link carries a spin $\tfrac12$, with Pauli
matrices $\sxb{j},\szb{j}$ acting on the $j$-th factor of $\LinkS_N:=(\C^2)^{\otimes N}$.  The
link $b_N$ is the \emph{boundary link}: it carries no hopping (there is no site $N+1$).  We set
$\sxb{0}:=\one$.  The total Hilbert space is
\begin{equation}\label{eq:HilbertN}
  \cH_N:=\cF_N\otimes\LinkS_N,\qquad \dim\cH_N=4^N ;
\end{equation}
the spins are bosonic, i.e.\ every $\sigma^{x,z}_{b_j}$ commutes with every $f_k$ (and every
spin operator is even).  We write $\sigma^z\ket{\uparrow}=\ket{\uparrow}$,
$\sigma^z\ket{\downarrow}=-\ket{\downarrow}$, $\ket{\pm}:=(\ket{\uparrow}\pm\ket{\downarrow})/\sqrt2$,
so that
\begin{equation}\label{eq:szflip}
  \sigma^x\ket{\pm}=\pm\ket{\pm},\qquad \sigma^z\ket{\pm}=\ket{\mp}\quad\text{(no phases).}
\end{equation}
For $x=(x_{b_1},\dots,x_{b_N})\in\{\pm1\}^N$ we write $\ket{x}_\sigma:=\bigotimes_{j}\ket{x_{b_j}}$
for the $\sigma^x$-eigenbasis (where $\ket{x_{b_j}}$ means $\ket{+}$ or $\ket{-}$), and for
$s\in\{\pm1\}^N$ we write $\ket{s}_z$ for the $\sigma^z$-eigenbasis; $\ket{\uparrow\cdots\uparrow}_z$
is the $z$-polarized state.  Link configurations are always written with subscripts $b_j$, so that
they are not confused with sites $x,y$.

\begin{definition}[The model]\label{def:model}
Let $J>0$ and $h\in\R\setminus\{0\}$.  The Hamiltonian on $\cH_N$ is
\begin{equation}\label{eq:H}
  H=-J\sum_{j=1}^{N-1}\szb{j}\bigl(f_j^\dagger f_{j+1}+f_{j+1}^\dagger f_j\bigr)
    -h\sum_{j=1}^{N}\sxb{j-1}\,\sxb{j}\qquad(\sxb{0}:=\one),
\end{equation}
and the Gauss operators are
\begin{equation}\label{eq:gauss}
  G_j:=\sxb{j-1}\,(-1)^{n_j}\,\sxb{j},\qquad j=1,\dots,N .
\end{equation}
An operator $O$ on $\cH_N$ is \emph{gauge invariant} if $[O,G_j]=0$ for all $j$.  For
$q\in\{\pm1\}^N$ let $P_q:=\prod_{j=1}^N\frac12(\one+q_jG_j)$.
\end{definition}

All algebraic statements of this section hold for arbitrary real $J,h$; the restrictions $J>0$,
$h\neq0$ only matter for the localization results.  (The simplicity of the spectrum of $\hq$, used in
\cref{thm:duality}(b) and \cref{rem:gaussproj}(ii), needs $J\neq0$.)  The boundary link enters $H$ only through
the term $-h\,\sxb{N-1}\sxb{N}$; for $N=1$, $H=-h\,\sxb{1}$.  Conjugation by $G_j$ is the local
gauge transformation at $j$: it maps $f_j\mapsto -f_j$, $\szb{j-1}\mapsto-\szb{j-1}$ (if $j\ge2$),
$\szb{j}\mapsto-\szb{j}$, and fixes all other $f_k$, $\sigma^z_{b_k}$ and all $\sigma^x_{b_k}$.

\begin{lemma}[Gauss law]\label{lem:gauss}
\leavevmode
\begin{enumerate}[label=(\roman*)]
\item The $G_j$ are self-adjoint unitaries, commute with each other, with $N_f$ and with $H$.
\item For every $j\in\{1,\dots,N\}$ one has the operator identity (one-dimensional $\Z_2$ Gauss law)
\begin{equation}\label{eq:gausslaw1d}
  \sxb{j}=(-1)^{N_{\le j}}\prod_{k\le j}G_k .
\end{equation}
In particular $\prod_{j=1}^NG_j=(-1)^{N_f}\sxb{N}$, and the boundary flux $\sxb{N}$ commutes with $H$.
\item For every $q\in\{\pm1\}^N$ the projection $P_q$ has rank $2^N$; thus $\cH_N=\bigoplus_q
P_q\cH_N$ with $2^N$ non-trivial superselection sectors, each invariant under $H$.
\end{enumerate}
\end{lemma}

\begin{proof}
(i) $G_j$ is a product of commuting self-adjoint unitaries ($\sigma^x$'s on distinct links and a
parity), and all $G_j$ are built from the mutually commuting operators $\sxb{k}$, $(-1)^{n_k}$.  They
commute with $N_f$ and with the electric term, which is a function of the $\sxb{k}$.  The hopping
term $T_k:=\szb{k}(f_k^\dagger f_{k+1}+f_{k+1}^\dagger f_k)$ anticommutes with $\sxb{k}$, with
$(-1)^{n_k}$ and with $(-1)^{n_{k+1}}$ (because $(-1)^{n_k}f_k^\dagger=-f_k^\dagger(-1)^{n_k}$ and
$(-1)^{n_k}$ commutes with $f_{k+1}$), and commutes with all other $\sxb{i}$ and $(-1)^{n_i}$.
$G_k$ contains exactly two of these anticommuting factors ($\sxb{k}$ and $(-1)^{n_k}$), as does
$G_{k+1}$ ($\sxb{k}$ and $(-1)^{n_{k+1}}$); the other $G_j$ contain none.  Hence $[G_j,T_k]=0$
for all $j,k$.
(ii) Since $(\sxb{k})^2=\one$, the product telescopes:
$\prod_{k\le j}G_k=\sxb{0}\sxb{j}\prod_{k\le j}(-1)^{n_k}=\sxb{j}(-1)^{N_{\le j}}$.  For $j=N$ this
gives $\prod_jG_j=(-1)^{N_f}\sxb{N}$, and $\sxb{N}=(-1)^{N_f}\prod_jG_j$ commutes with $H$ by (i).
(iii) The vectors $\ket{n}\otimes\ket{x}_\sigma$ form an orthonormal basis of joint eigenvectors of
the $G_j$, with eigenvalues $x_{b_{j-1}}(-1)^{n_j}x_{b_j}$ ($x_{b_0}:=1$).  For fixed $n$ the map
$x\mapsto q$, $q_j=x_{b_{j-1}}(-1)^{n_j}x_{b_j}$, is a bijection of $\{\pm1\}^N$ (its inverse is
$x_{b_j}=\prod_{k\le j}q_k(-1)^{n_k}$).  Hence every $q$ is attained by exactly $2^N$ basis vectors.
\end{proof}

The eigenvalue $q_j$ of $G_j$ is the static $\Z_2$ background charge at site $j$.  The
\emph{gauge-invariant subspace} in the usual lattice-gauge-theory sense (no background charges) is
$P_{(1,\dots,1)}\cH_N$; the $z$-polarized initial states of \cite{SKKM2017} are \emph{not} in it,
but are uniform superpositions over all sectors (\cref{cor:sector}).

\begin{remark}[Comparison with \cite{SKKM2017}; boundary conventions]\label{rem:boundary}
\leavevmode
\begin{enumerate}[label=(\alph*)]
\item \emph{Hamiltonian.}  \eqref{eq:H} is Eq.~(1) of \cite{SKKM2017},
$\hat H=-J\sum_{\langle ij\rangle}\hat\sigma^z_{i,j}\hat f_i^\dagger\hat f_j
-h\sum_i\hat\sigma^x_{i-1,i}\hat\sigma^x_{i,i+1}$ (the sum over $\langle ij\rangle$ contains both
orientations).  The electric term is the \emph{two-link} term; this is essential.  The standard
$\Z_2$ gauge theory with single-link electric term $-h\sum_b\sigma^x_b$ is \emph{not} free:
already one fermion sees a random-walk potential and two fermions interact
(\cref{rem:randomwalk}).
\item \emph{Duality.}  Following the Kramers--Wannier duality of the Ising chain
\cite{KramersWannier1941,FradkinSusskind1978}, SKKM introduce site spins by
$\hat\tau^z_j=\hat\sigma^x_{j-1,j}\hat\sigma^x_{j,j+1}$,
$\hat\sigma^z_{j,j+1}=\hat\tau^x_j\hat\tau^x_{j+1}$ (their Eq.~(3)), the charges
$\hat q_j=\hat\tau^z_j(-1)^{\hat n_j}$, which are our $G_j$, and the fermions
$\hat c_j=\hat\tau^x_j\hat f_j$, which lead to their Eqs.~(4) and (2),
$\hat H_{\{q_j\}}=-J\sum_{\langle ij\rangle}\hat c_i^\dagger\hat c_j+2h\sum_jq_j(\hat c_j^\dagger\hat c_j-\tfrac12)$.
Since $2h\sum_jq_j(n_j-\frac12)=2h\sum_jq_jn_j-h\sum_jq_j$, this is exactly
$\dGamma(\hq)-h\sum_jq_j$ in \cref{thm:duality}(b).  With the boundary link the operators
$\hat\tau^z_j:=\sxb{j-1}\sxb{j}$ and $\hat\tau^x_j:=\prod_{k\ge j}\szb{k}$ satisfy the Pauli
algebra exactly on $\LinkS_N$, and $\hat\tau^x_j\hat\tau^x_{j+1}=\szb{j}$ for $j<N$; with this
choice SKKM's $\hat c_j$ is our dressed fermion $\cdr_j$ of \eqref{eq:dressed} and SKKM's $\hat\tau^x_j$
is $U\tauz{j}U^\dagger$.  (Our dual spins $\tau_j$ of \cref{thm:duality} live on a different space
and should not be confused with SKKM's $\hat\tau_j$: the charge is $U^\dagger G_jU=\taux{j}$.)
\item \emph{Initial states.}  For $\ket{S}=\ket{\uparrow\uparrow\cdots}$ SKKM write the initial state as
\[
  \ket{0}=2^{-N/2}\sum_{\{q_i\}=\pm1}\ket{q_1,\dots,q_N}\otimes\ket{\psi},
\]
a sum over all $2^N$ charge configurations (their Eq.~(5)).  With the boundary link this holds literally
(\cref{cor:sector}(a)).  SKKM work with $N$ sites and open boundary conditions, and they
consider domain-wall, density-wave and $h=0$ ground-state fermion states, with $N$ up to $200$.
\item \emph{The ``open'' chain} has $N$ sites and $N-1$ links, $\sxb{0}=\sxb{N}:=\one$ in
\eqref{eq:H} and \eqref{eq:gauss}.  Then $\prod_jG_j=(-1)^{N_f}$ identically, so only the
charge configurations with $\prod_jq_j=(-1)^{N_f}$ occur, and for a matter state of definite
parity the sum in SKKM's Eq.~(5) runs over the $2^{N-1}$ configurations of a parity hyperplane.
The open chain is the restriction of our model to the invariant subspace $\{\sxb{N}=+1\}$
(\cref{prop:variants}(i)).  The boundary link $b_N$ is a static $\Z_2$ flux through the right
end; it makes $\sxb{N}=(-1)^{N_f}\prod_jG_j$ a free conserved label, removes the parity constraint,
and makes the charges exactly i.i.d.  For the $z$-polarized gauge state and \emph{any} matter state, the
sector laws of the two conventions differ by a density bounded by $2$ (\cref{prop:variants}(i)).  This is
harmless for all localization \emph{upper} bounds (\cref{rem:productlaws}(c)).
\item \emph{The ``ring''} ($N$ sites, $N$ links, hopping also across $b_N=(N,1)$) has, in
addition, the conserved Wilson loop $W=\prod_j\szb{j}$, again the constraint
$\prod_jG_j=(-1)^{N_f}$, and the one-body operator acquires the twist $W$ on the bond $(N,1)$,
with no additional parity-dependent twist (\cref{prop:variants}(ii)).  SKKM use periodic boundary
conditions for their Fermi-sea initial state and note the global constraint, which is satisfied by
the $z$-polarized spin state; in our language this state has $W=+1$.
\end{enumerate}
\end{remark}

\subsection{The duality}\label{sec:model:duality}

\paragraph{Dual space.}
Let $\cK_N:=\cF^c_N\otimes\SiteT_N$, where $\cF^c_N$ is the Fock space of fermions $c_1,\dots,c_N$
with occupation basis $\ket{n}_c$ defined by the same ordering convention \eqref{eq:occbasis}, and
$\SiteT_N=(\C^2)^{\otimes N}$ carries spins $\tau_j$ on the \emph{sites}, commuting with the $c$'s.
We use the orthonormal basis $\ket{q}$, $q\in\{\pm1\}^N$, of $\SiteT_N$ with
\begin{equation}\label{eq:taubasis}
  \taux{j}\ket{q}=q_j\ket{q},\qquad \tauz{j}\ket{q}=\ket{q^{(j)}},
\end{equation}
where $q^{(j)}$ is $q$ with $q_j$ replaced by $-q_j$ (this fixes the phases of the basis).  We
write $\ket{n,q}:=\ket{n}_c\otimes\ket{q}$, set $\tauz{N+1}:=\one$, and identify $\cF_N$ with
$\cF_N^c$ via $\ket{n}\mapsto\ket{n}_c$; under this identification $f_j$ corresponds to $c_j$, and
a matter vector $\psi\in\cF_N$ is also regarded as a vector of $\cF_N^c$.  For an $N\times N$
matrix $A$ let $\dGamma(A):=\sum_{x,y}A_{xy}c_x^\dagger c_y$, and for a unitary $V$ let $\Gamma(V)$
be its second quantization ($\Gamma(V)\ket{0}=\ket{0}$, $\Gamma(V)c^\dagger(u)\Gamma(V)^*=c^\dagger(Vu)$),
so that $\e^{-\ii t\dGamma(A)}=\Gamma(\e^{-\ii tA})$ for self-adjoint $A$.  Let $\Delta_N$ be the
adjacency matrix of the path $\{1,\dots,N\}$, $(\Delta_Nu)(x)=u(x+1)+u(x-1)$ with
$u(0)=u(N+1)=0$, and
\begin{equation}\label{eq:hq}
  \hq:=-J\Delta_N+2h\,\diag(q)\quad\text{on }\ell^2(\{1,\dots,N\}),\qquad \lambda:=2h/J .
\end{equation}

\paragraph{The unitary.}
For $n\in\{0,1\}^N$ and $q\in\{\pm1\}^N$ define the link configuration
\begin{equation}\label{eq:xnq}
  x_{b_j}(n,q):=\prod_{k\le j}q_k(-1)^{n_k},\qquad j=1,\dots,N,
\end{equation}
and the linear map $U:\cK_N\to\cH_N$ by
\begin{equation}\label{eq:U}
  U\ket{n,q}:=\ket{n}\otimes\ket{x(n,q)}_\sigma .
\end{equation}
Conversely, with $x_{b_0}:=1$,
\begin{equation}\label{eq:qofx}
  q_j=x_{b_{j-1}}\,(-1)^{n_j}\,x_{b_j},\qquad j=1,\dots,N .
\end{equation}

\begin{theorem}[Exact duality]\label{thm:duality}
$U$ is unitary, preserves all occupation numbers (in particular it is even), and:
\begin{enumerate}[label=(\alph*)]
\item \emph{Gauss operators:} $U^\dagger G_jU=\taux{j}$ for all $j$; hence
$U^\dagger P_qU=\one\otimes\ket{q}\bra{q}$.
\item \emph{Hamiltonian:}
\begin{equation}\label{eq:dualH}
  U^\dagger HU=-J\sum_{j=1}^{N-1}(c_j^\dagger c_{j+1}+c_{j+1}^\dagger c_j)-h\sum_{j=1}^N\taux{j}(-1)^{n_j}
  =\sum_{q\in\{\pm1\}^N}\Bigl(\dGamma(\hq)-h\sum_{j=1}^Nq_j\Bigr)\otimes\ket{q}\bra{q}.
\end{equation}
Consequently, with $E_1(q)<\dots<E_N(q)$ the eigenvalues of $\hq$ (simple, since $J\neq0$ and $\hq$ is a
Jacobi matrix with non-zero off-diagonal entries), the spectrum of $H$ (with
multiplicities) is the union over $q$ of the multisets
$\{\sum_\alpha m_\alpha E_\alpha(q)-h\sum_jq_j:\ m\in\{0,1\}^N\}$.
\item \emph{Dictionary:} for all $j$ (with $\tauz{N+1}=\one$),
\begin{gather*}
  U^\dagger n_jU=n_j,\qquad U^\dagger(-1)^{N_f}U=(-1)^{N_c},\qquad U^\dagger f_jU=c_j\tauz{j},\\
  U^\dagger\szb{j}U=\tauz{j}\tauz{j+1},\qquad
  U^\dagger\Bigl(\prod_{k=j}^N\szb{k}\Bigr)U=\tauz{j},\\
  U^\dagger\sxb{j}U=\prod_{k\le j}\taux{k}(-1)^{n_k},\qquad
  U^\dagger\sxb{j-1}\sxb{j}U=\taux{j}(-1)^{n_j},\\
  U^\dagger\bigl(\szb{j}f_j^\dagger f_{j+1}\bigr)U=c_j^\dagger c_{j+1}\quad(j<N).
\end{gather*}
\item \emph{Dressed fermions:} the operators
\begin{equation}\label{eq:dressed}
  \cdr_j:=Uc_jU^\dagger=f_j\prod_{k=j}^{N}\szb{k}
\end{equation}
satisfy the CAR, commute with every $G_k$, are odd, and $\cdr_j^\dagger\cdr_j=n_j$.  For $x<y$,
$\cdr_x^\dagger\cdr_y=f_x^\dagger\bigl(\prod_{x\le k<y}\szb{k}\bigr)f_y$, the gauge-invariant
hopping with a Wilson line between $x$ and $y$.  In terms of the original operators,
\begin{equation}\label{eq:Hdressed}
  H=-J\sum_{j=1}^{N-1}(\cdr_j^\dagger\cdr_{j+1}+\cdr_{j+1}^\dagger\cdr_j)-h\sum_{j=1}^NG_j(-1)^{n_j}
   =\sum_q P_q\Bigl(\sum_{x,y}(\hq)_{xy}\cdr_x^\dagger\cdr_y-h\sum_jq_j\Bigr).
\end{equation}
\item \emph{Dynamics:}
$U^\dagger\e^{-\ii tH}U=\sum_q\e^{\ii ht\sum_jq_j}\,\Gamma(\e^{-\ii t\hq})\otimes\ket{q}\bra{q}$, and
\begin{equation}\label{eq:heisenberg}
  \e^{\ii tH}\cdr_x\e^{-\ii tH}=\sum_qP_q\sum_{y=1}^N(\e^{-\ii t\hq})_{xy}\,\cdr_y\qquad(t\in\R).
\end{equation}
\end{enumerate}
\end{theorem}

\begin{proof}
\emph{Unitarity.}  By \eqref{eq:xnq}--\eqref{eq:qofx} the map $(n,q)\mapsto(n,x(n,q))$ is a
bijection of $\{0,1\}^N\times\{\pm1\}^N$.  Hence $U$ maps the orthonormal basis
$\{\ket{n,q}\}$ of $\cK_N$ onto the orthonormal basis $\{\ket{n}\otimes\ket{x}_\sigma\}$ of
$\cH_N$, and it does not change $n$.

\emph{(a).}  $G_j(\ket{n}\otimes\ket{x}_\sigma)=x_{b_{j-1}}(-1)^{n_j}x_{b_j}\ket{n}\otimes\ket{x}_\sigma$.
For $x=x(n,q)$ the eigenvalue is $q_j$ by \eqref{eq:qofx}, so $G_jU\ket{n,q}=q_jU\ket{n,q}=U\taux{j}\ket{n,q}$.
The formula for $P_q$ follows.

\emph{(c).}  The first two identities hold because $U$ does not change $n$.
$\sxb{j}U\ket{n,q}=x_{b_j}(n,q)U\ket{n,q}$, and $x_{b_j}(n,q)$ is the eigenvalue of
$\prod_{k\le j}\taux{k}(-1)^{n_k}$ on $\ket{n,q}$; the identity for $\sxb{j-1}\sxb{j}$ follows by
multiplying the cases $j-1$ and $j$.  By \eqref{eq:szflip}, $\szb{j}\ket{x}_\sigma=\ket{x'}_\sigma$
where $x'$ is $x$ with $x_{b_j}$ flipped.  By \eqref{eq:qofx}, $x_{b_j}$ enters only $q_j$ and
$q_{j+1}$ (only $q_N$ if $j=N$), so $x(n,q)'=x(n,q')$ with $q'$ obtained from $q$ by flipping
$q_j$ and $q_{j+1}$ (resp.\ $q_N$); thus $\szb{j}U\ket{n,q}=U\tauz{j}\tauz{j+1}\ket{n,q}$.  For the
fermions, \eqref{eq:fsign} gives
$f_jU\ket{n,q}=\delta_{n_j,1}(-1)^{N_{<j}(n)}\ket{n-e_j}\otimes\ket{x(n,q)}_\sigma$.  Now
$x(n-e_j,q^{(j)})=x(n,q)$: in \eqref{eq:xnq} only the factor $q_j(-1)^{n_j}$ changes, into
$(-q_j)(-(-1)^{n_j})$.  Hence, using \eqref{eq:fsign} for the $c$'s,
\[
  f_jU\ket{n,q}=\delta_{n_j,1}(-1)^{N_{<j}(n)}\,U\ket{n-e_j,q^{(j)}}=U\,c_j\tauz{j}\ket{n,q}.
\]
The fermionic sign $(-1)^{N_{<j}(n)}$ is the same on both sides because $U$ does not change the
occupation numbers; this is the only place where the ordering convention enters, and the result
holds for any convention used simultaneously for $f$ and $c$.  Since the $\tau$'s commute with the
$c$'s and $(\tauz{k})^2=\one$,
$U^\dagger\szb{j}f_j^\dagger f_{j+1}U=\tauz{j}\tauz{j+1}\,\tauz{j}c_j^\dagger\,c_{j+1}\tauz{j+1}=c_j^\dagger c_{j+1}$.
Finally $\prod_{k=j}^N\tauz{k}\tauz{k+1}=\tauz{j}$ (telescoping, $\tauz{N+1}=\one$).

\emph{(b).}  By (c), $(\szb{j}f_j^\dagger f_{j+1})^\dagger=\szb{j}f_{j+1}^\dagger f_j\mapsto c_{j+1}^\dagger c_j$
and $\sxb{j-1}\sxb{j}\mapsto\taux{j}(-1)^{n_j}$, which gives the first expression in
\eqref{eq:dualH}.  On $\cF^c_N\otimes\ket{q}$ the electric part acts as
$-h\sum_jq_j(\one-2n_j)=2h\sum_jq_jn_j-h\sum_jq_j$, and
$-J\sum_{j<N}(c_j^\dagger c_{j+1}+c_{j+1}^\dagger c_j)+2h\sum_jq_jn_j=\dGamma(\hq)$.  If
$\hq\varphi_\alpha=E_\alpha\varphi_\alpha$ with $(\varphi_\alpha)$ orthonormal and
$a_\alpha:=\sum_x\overline{\varphi_\alpha(x)}c_x$, then $\dGamma(\hq)=\sum_\alpha E_\alpha a_\alpha^\dagger a_\alpha$,
whose eigenvalues are $\sum_\alpha m_\alpha E_\alpha$, $m\in\{0,1\}^N$.

\emph{(d).}  By (c), $c_j=c_j\tauz{j}\tauz{j}=(U^\dagger f_jU)\bigl(U^\dagger\prod_{k\ge j}\szb{k}U\bigr)$;
conjugating with $U$ gives \eqref{eq:dressed}.  The CAR and oddness are preserved by the even unitary
$U$; $c_j$ commutes with $\taux{k}$, so $\cdr_j$ commutes with $G_k$ by (a); and
$\cdr_j^\dagger\cdr_j=Un_jU^\dagger=n_j$ by (c).  For $x<y$ the strings
$\prod_{k\ge x}\szb{k}\prod_{k\ge y}\szb{k}=\prod_{x\le k<y}\szb{k}$ combine and commute with the
$f$'s.  \eqref{eq:Hdressed} is $U(\eqref{eq:dualH})U^\dagger$, using
$U\taux{j}(-1)^{n_j}U^\dagger=G_j(-1)^{n_j}$ and $U(\one\otimes\ket{q}\bra{q})U^\dagger=P_q$.

\emph{(e).}  The first identity follows from \eqref{eq:dualH}.  For a self-adjoint $A$, both sides of
$\e^{\ii t\dGamma(A)}c_x\e^{-\ii t\dGamma(A)}=\sum_y(\e^{-\ii tA})_{xy}c_y$ solve the same linear
differential equation with the same initial value, because $[c_a^\dagger c_b,c_x]=-\delta_{ax}c_b$
gives $\ii[\dGamma(A),c_x]=-\ii\sum_bA_{xb}c_b$.  Conjugating with $U$, the scalar phases
$\e^{\pm\ii ht\sum q}$ cancel since $c_y$ commutes with $\ket{q}\bra{q}$, and
$U(c_y\otimes\ket{q}\bra{q})U^\dagger=\cdr_yP_q$.
\end{proof}

\begin{remark}[Where the Wilson line ends]\label{rem:strings}
The string of $\cdr_j$ runs from $b_j$ to the boundary link $b_N$.  A left-anchored string does not
work: $c'_j:=(\prod_{k<j}\szb{k})f_j=\cdr_j\prod_{k=1}^N\szb{k}$ satisfies the CAR but anticommutes
with $G_1$ (for $j=1$ through $(-1)^{n_1}$, for $j\ge2$ through $\szb{1}$), because
$U^\dagger\prod_k\szb{k}U=\tauz{1}$ flips $q_1$.  In the open and ring variants of
\cref{rem:boundary}, $\prod_jG_j=(-1)^{N_f}$, so every gauge-invariant operator is even and there is
no gauge-invariant single fermion at all; only bilinears such as
$f_x^\dagger(\prod_{x\le k<y}\szb{k})f_y$ survive.  In our model $\cdr_j$ anticommutes with both
$(-1)^{N_f}$ and $\sxb{N}$, consistently with $\prod_jG_j=(-1)^{N_f}\sxb{N}$.  The algebra generated
by $\{\cdr_j,G_j,\prod_{k\ge j}\szb{k}\}_{j\le N}$ is a CAR$\,\otimes\,$Pauli system on $\cH_N$, and $U$
is the unitary that intertwines it with $\{c_j,\taux{j},\tauz{j}\}$.
\end{remark}

\begin{proposition}[Open and ring variants]\label{prop:variants}
\leavevmode
\begin{enumerate}[label=(\roman*)]
\item \emph{Open chain.}  Let $H^{\rm open}$ and $G^{\rm open}_j$ be given by \eqref{eq:H},
\eqref{eq:gauss} on $\cF_N\otimes(\C^2)^{\otimes(N-1)}$ with $\sxb{0}=\sxb{N}:=\one$, and let
$\iota\chi:=\chi\otimes\ket{+}_{b_N}$.  Then $H(\one\otimes\iota)=(\one\otimes\iota)H^{\rm open}$ and
$G_j(\one\otimes\iota)=(\one\otimes\iota)G_j^{\rm open}$, so $H^{\rm open}$ is unitarily equivalent to
$H$ restricted to the invariant subspace $\{\sxb{N}=+1\}$, and
\begin{equation}\label{eq:openparity}
  U^\dagger\tfrac12(\one+\sxb{N})U=\sum_q\Pi_{\pi(q)}\otimes\ket{q}\bra{q},\qquad\pi(q):=\prod_{j}q_j,
\end{equation}
where $\Pi_{\pm}$ is the projection onto $\{(-1)^{N_c}=\pm1\}$ in $\cF^c_N$.  Thus
$H^{\rm open}\cong\bigoplus_q\bigl(\dGamma(\hq)-h\sum_jq_j\bigr)\big|_{\Pi_{\pi(q)}\cF^c_N}$ (the
parity constraint $\prod_jq_j=(-1)^{N_f}$).  For every unit vector $\psi\in\cF_N$, the state
$\psi\otimes\ket{\uparrow\cdots\uparrow}_z\in\cF_N\otimes(\C^2)^{\otimes(N-1)}$ satisfies
\begin{equation}\label{eq:openinit}
  U^\dagger(\one\otimes\iota)\bigl(\psi\otimes\ket{\uparrow\cdots\uparrow}_z\bigr)
  =2^{-(N-1)/2}\sum_q\Pi_{\pi(q)}\psi\otimes\ket{q}.
\end{equation}
So its sector law $\mu(q)=2^{-(N-1)}\norm{\Pi_{\pi(q)}\psi}^2$ has density
$2\norm{\Pi_{\pi(q)}\psi}^2\le2$ with respect to the i.i.d.\ Bernoulli$(1/2)$ law, and the sector $q$
carries the matter state $\Pi_{\pi(q)}\psi$ (normalized).  If $\psi$ has definite parity $\pi$, the sector
law is uniform on $\{q:\prod_jq_j=\pi\}$, with density $2\cdot\one[\prod_jq_j=\pi]$, and every sector in
its support carries the same matter state $\psi$.
\item \emph{Ring} ($N\ge3$).  Let $H^{\rm ring}$ be \eqref{eq:H} on $\cH_N$ with the hopping sum
extended to $j=N$ ($f_{N+1}:=f_1$) and $\sxb{0}:=\sxb{N}$, let $G^{\rm ring}_j$ be \eqref{eq:gauss}
with $\sxb{0}:=\sxb{N}$, and $W:=\prod_{j=1}^N\szb{j}$.  Then, with the same $U$,
\begin{gather*}
  U^\dagger G^{\rm ring}_1U=(-1)^{N_c}\prod_{k\ge2}\taux{k},\qquad U^\dagger G^{\rm ring}_jU=\taux{j}\ (j\ge2),
  \qquad U^\dagger WU=\tauz{1},\\
  \begin{aligned}
  U^\dagger H^{\rm ring}U={}&-J\sum_{j=1}^{N-1}(c_j^\dagger c_{j+1}+{\rm h.c.})-J\tauz{1}(c_N^\dagger c_1+c_1^\dagger c_N)\\
  &-h(-1)^{N_c}\Bigl(\prod_{k\ge2}\taux{k}\Bigr)(-1)^{n_1}-h\sum_{j=2}^N\taux{j}(-1)^{n_j}.
  \end{aligned}
\end{gather*}
Consequently $W$ and the $G^{\rm ring}_j$ commute with $H^{\rm ring}$, $\prod_jG_j^{\rm ring}=(-1)^{N_f}$, and on
the joint eigenspace $\{G^{\rm ring}_j=q_j\ \forall j,\ W=w\}$ (non-zero for every $q$ and $w$) the
Hamiltonian $H^{\rm ring}$ is unitarily equivalent to $\dGamma(h_{q,w})-h\sum_jq_j$ restricted to
$\Pi_{\pi(q)}\cF^c_N$, where
\begin{equation}\label{eq:ringtwist}
  h_{q,w}:=\hq-Jw\bigl(\ket{1}\bra{N}+\ket{N}\bra{1}\bigr).
\end{equation}
The twist on the bond $(N,1)$ is exactly $w$, the eigenvalue of the Wilson loop, with no
parity-dependent factor.  The state $\psi\otimes\ket{\uparrow\cdots\uparrow}_z$ with $\psi$ of parity
$\pi$ has $W=+1$ (periodic $h_{q,1}$) and sector law uniform on $\{q:\prod_jq_j=\pi\}$.
\end{enumerate}
\end{proposition}

\begin{proof}
(i) $\sxb{N}$ enters $H$ only through $-h\sxb{N-1}\sxb{N}$ and $G_j$ only for $j=N$; on
$\ket{+}_{b_N}$ it acts as $1$, which gives the intertwining relations.  By
\cref{lem:gauss}(ii) and \cref{thm:duality}(a),(c),
$U^\dagger\sxb{N}U=(-1)^{N_c}\prod_j\taux{j}$, which is \eqref{eq:openparity}.  Restricting
\eqref{eq:dualH} to its range gives the sector decomposition.  For the initial state,
$\iota(\ket{\uparrow\cdots\uparrow}_z)=2^{-(N-1)/2}\sum_{x:\,x_{b_N}=+1}\ket{x}_\sigma$.  For a Fock
vector $\ket{n}$, by \eqref{eq:U} and \eqref{eq:qofx} the vectors $U^\dagger(\ket{n}\otimes\ket{x}_\sigma)=\ket{n,q(n,x)}$
with $x_{b_N}=+1$ are exactly the $\ket{n,q}$ with $\prod_jq_j=(-1)^{|n|}$, each once (indeed
$\prod_jq_j=x_{b_0}x_{b_N}(-1)^{|n|}$ by \eqref{eq:qofx}).  By linearity, for $\psi=\sum_n\psi(n)\ket{n}$,
\[
  U^\dagger(\psi\otimes\iota\ket{\uparrow\cdots\uparrow}_z)
  =2^{-(N-1)/2}\sum_q\Bigl(\sum_{n:\,(-1)^{|n|}=\pi(q)}\psi(n)\ket{n}\Bigr)\otimes\ket{q}
  =2^{-(N-1)/2}\sum_q\Pi_{\pi(q)}\psi\otimes\ket{q},
\]
which is \eqref{eq:openinit}; the statements on the sector law follow.

(ii) The operators $f_j,\sigma_{b_j}$ and $U$ are the same as before, so the dictionary of
\cref{thm:duality}(c) applies.  The new hopping term maps as
$\szb{N}f_N^\dagger f_1\mapsto\tauz{N}\,c_N^\dagger\tauz{N}\,c_1\tauz{1}=\tauz{1}c_N^\dagger c_1$, and similarly
for its adjoint.  The new electric term maps as
\[
  \sxb{N}\sxb{1}\mapsto\Bigl(\prod_{k\le N}\taux{k}(-1)^{n_k}\Bigr)\taux{1}(-1)^{n_1}
  =(-1)^{N_c}\Bigl(\prod_{k\ge2}\taux{k}\Bigr)(-1)^{n_1};
\]
multiplying by $(-1)^{n_1}$ gives $U^\dagger G^{\rm ring}_1U$.  $W\mapsto\prod_j\tauz{j}\tauz{j+1}=\tauz{1}$.  The dual Hamiltonian
contains $\tau_1$ only through $\tauz{1}$ and $\tau_k$, $k\ge2$, only through $\taux{k}$, and it
commutes with $(-1)^{N_c}$.  Use the basis $\ket{w}_1\otimes\ket{q_2,\dots,q_N}$ of $\SiteT_N$ with
$\tauz{1}\ket{w}_1=w\ket{w}_1$.  On $\Pi_\pi\cF_N^c\otimes\ket{w}_1\otimes\ket{q_2,\dots,q_N}$ put
$q_1:=\pi\prod_{k\ge2}q_k$.  Then $U^\dagger G^{\rm ring}_1U=q_1$ there, and the dual Hamiltonian acts as
\[
  -J\sum_{j<N}(c_j^\dagger c_{j+1}+{\rm h.c.})-Jw(c_N^\dagger c_1+{\rm h.c.})-h\sum_jq_j(\one-2n_j)
  =\dGamma(h_{q,w})-h\sum_jq_j .
\]  These subspaces exhaust $\cK_N$ and are the joint eigenspaces of
$(U^\dagger G_j^{\rm ring}U)_j$ and $U^\dagger WU$, with $\pi=\prod_jq_j$.  Finally
$U^\dagger(\psi\otimes\ket{\uparrow\cdots\uparrow}_z)=2^{-N/2}\sum_q\psi\otimes\ket{q}$ (the same
computation as in \cref{cor:sector}(a) below), and
$2^{-1/2}\sum_{q_1}\ket{q_1}$ is the $\tauz{1}=+1$ eigenvector by \eqref{eq:taubasis}; so $w=1$.  Since
$U^\dagger G_1^{\rm ring}U=\pi\prod_{k\ge2}\taux{k}$ on $\Pi_\pi\cF^c_N\otimes\SiteT_N$, the law of
$(q_2,\dots,q_N)$ is uniform and $q_1=\pi\prod_{k\ge2}q_k$.
\end{proof}

\subsection{Locality of the duality across a cut}\label{sec:model:cut}

Fix $1\le\ell<N$.  Let $A$ consist of the sites $1,\dots,\ell$ and the links $b_1,\dots,b_\ell$, and
$B$ of the sites $\ell+1,\dots,N$ and the links $b_{\ell+1},\dots,b_N$.  Identify
$\cF_N\cong\cF_{\le\ell}\otimes\cF_{>\ell}$ via $\ket{n}\mapsto\ket{n_{\le\ell}}\otimes\ket{n_{>\ell}}$, where
$\ket{n_{\le\ell}}=(f_1^\dagger)^{n_1}\cdots(f_\ell^\dagger)^{n_\ell}\ket0$ and
$\ket{n_{>\ell}}=(f_{\ell+1}^\dagger)^{n_{\ell+1}}\cdots(f_N^\dagger)^{n_N}\ket0$.  By
\eqref{eq:fsign}, $f_j\mapsto f_j\otimes\one$ for $j\le\ell$ and $f_j\mapsto(-1)^{N_{\le\ell}}\otimes f_j$
for $j>\ell$.  Hence the CAR algebra of the modes $1,\dots,\ell$ is exactly $\mathcal{B}(\cF_{\le\ell})\otimes\one$,
and
\[
  \cH_N\cong\cH_A\otimes\cH_B,\qquad\cH_A:=\cF_{\le\ell}\otimes(\C^2)^{\otimes\ell},\quad
  \cH_B:=\cF_{>\ell}\otimes(\C^2)^{\otimes(N-\ell)},
\]
with the algebra generated by $\{f_j,\sigma^{x,z}_{b_j}\}_{j\le\ell}$ equal to $\mathcal{B}(\cH_A)\otimes\one$.
For a unit vector $\Psi\in\cH_N$ we let $\rho_A(\Psi):=\Tr_{\cH_B}\ket{\Psi}\bra{\Psi}$ and
$S_A(\Psi):=-\Tr\rho_A\ln\rho_A$.  Thus $\Tr(\rho_A(\Psi)O)=\langle\Psi,O\Psi\rangle$ for every $O$ in the
algebra of $A$, so $\rho_A$ is the canonical fermionic reduced state of the left segment, and
$\Sfull(t)=S_A(\Psi(t))$.  On the dual side we use the same identification
$\cK_N\cong\cK_{A'}\otimes\cK_{B'}$, $\ket{n,q}\mapsto\ket{n_{\le\ell},q_{\le\ell}}\otimes\ket{n_{>\ell},q_{>\ell}}$, where
$A'=(c_k,\tau_k)_{k\le\ell}$, and we write $\rho_{A'}(\Phi)$, $S_{A'}(\Phi)$.

Let $U_{\le\ell}:\cK_{A'}\to\cH_A$ be the duality unitary \eqref{eq:U} of the chain $\{1,\dots,\ell\}$
(whose boundary link is $b_\ell$), and $U_{>\ell}:\cK_{B'}\to\cH_B$,
$U_{>\ell}\ket{n_{>\ell},q_{>\ell}}=\ket{n_{>\ell}}\otimes\ket{y}_\sigma$ with
$y_{b_j}:=\prod_{\ell<k\le j}q_k(-1)^{n_k}$, $j>\ell$.  Let $P^\pm_\ell:=\frac12(\one\pm\sxb{\ell})$ on
$\cH_A$, $V_\ell:=\prod_{k=\ell+1}^N\szb{k}$ on $\cH_B$, and
\begin{equation}\label{eq:Cgate}
  C_\ell:=P^+_\ell\otimes\one+P^-_\ell\otimes V_\ell .
\end{equation}

\begin{lemma}[Cut-locality]\label{lem:cutlocal}
\leavevmode
\begin{enumerate}[label=(\alph*)]
\item $U=C_\ell\,(U_{\le\ell}\otimes U_{>\ell})$.
\item For every unit vector $\Phi\in\cK_N$: $\ \abs{S_A(U\Phi)-S_{A'}(\Phi)}\le\ln2$.
\item If $\Phi$ is an eigenvector of $(-1)^{N_c}$ and of $\prod_{k\le\ell}\taux{k}$ (which is
$U^\dagger\prod_{k\le\ell}G_kU$, the total $\Z_2$ charge in $A$), then $S_A(U\Phi)=S_{A'}(\Phi)$.  In
particular, for a definite-charge state $\Phi=\phi\otimes\ket{q}$ with $\phi\in\cF_N^c$ of definite
parity, $S_A(U\Phi)=S(\rho^\phi_{\le\ell})$, the entanglement entropy of $\phi$ for the modes
$1,\dots,\ell$.
\item The constant $\ln2$ in (b) is attained in both directions (already for $N=2$, $\ell=1$).
\end{enumerate}
\end{lemma}

In particular, for a definite-charge initial state $\Psi_0=\ket{\xi}_\sigma\otimes\ket{n}$
(\cref{lem:productlaw}(d)) with $q=q(n,\xi)$, \cref{thm:duality}(e) gives
$U^\dagger\Psi(t)=\e^{\ii ht\sum_jq_j}\psi_q(t)\otimes\ket{q}$, where $\psi_q(t)=\Gamma(\e^{-\ii t\hq})\ket{n}$
has parity $(-1)^{\abs{n}}$.  Hence, for all $t$,
\begin{equation}\label{eq:quenchedSfull}
  \Sfull(t)=S_A(\Psi(t))=S\bigl(\rho^{\psi_q(t)}_{\le\ell}\bigr):
\end{equation}
the full matter-plus-gauge entanglement of a definite-charge state is exactly the free-fermion
entanglement in its sector.

The defect in (b) comes only from the electric flux through the cut link $b_\ell$.  By
\eqref{eq:gausslaw1d}, $\sxb{\ell}=(-1)^{N_{\le\ell}}\prod_{k\le\ell}G_k$, i.e.\ the flux is fixed by
the matter parity and the background charge of $A$.  It is a single classical bit, which explains
the bound $\ln2$.

\begin{lemma}[Two-block pinching]\label{lem:pinching}
Let $P^++P^-=\one$ be orthogonal projections on a finite-dimensional space and
$\mathcal{P}(\sigma):=P^+\sigma P^++P^-\sigma P^-$.  Then $S(\sigma)\le S(\mathcal{P}(\sigma))\le S(\sigma)+\ln2$
for every density matrix $\sigma$.
\end{lemma}

This is standard (a pinching into $k$ blocks raises the entropy by at most $\ln k$); we include the
short proof.

\begin{proof}
With the unitary $R:=P^+-P^-$ one has $\mathcal{P}(\sigma)=\frac12(\sigma+R\sigma R^\dagger)$.  The lower
bound is concavity of $S$, since $S(R\sigma R^\dagger)=S(\sigma)$.  For the upper bound let
$\sigma_1,\sigma_2$ be density matrices and $\bar\sigma=\frac12(\sigma_1+\sigma_2)$.  For $\varepsilon>0$,
$\bar\sigma+\varepsilon\ge\frac12\sigma_i+\varepsilon>0$, and operator monotonicity of the logarithm gives
$-\Tr\sigma_i\ln(\bar\sigma+\varepsilon)\le-\Tr\sigma_i\ln(\tfrac12\sigma_i+\varepsilon)$.  Averaging over
$i=1,2$ and letting $\varepsilon\downarrow0$ (with $0\ln0=0$) yields
$S(\bar\sigma)\le\frac12\sum_iS(\sigma_i)+\ln2$.  Apply this with $\sigma_1=\sigma$, $\sigma_2=R\sigma R^\dagger$.
\end{proof}

\begin{proof}[Proof of \cref{lem:cutlocal}]
(a) Write $n=(n_<,n_>)$ and $q=(q_<,q_>)$ for the parts in $A'$ and $B'$.  By \eqref{eq:xnq},
$x_{b_j}(n,q)=x_{b_j}(n_<,q_<)$ for $j\le\ell$ and $x_{b_j}(n,q)=x_{b_\ell}(n_<,q_<)\,y_{b_j}(n_>,q_>)$
for $j>\ell$.  Under the identifications, $(U_{\le\ell}\otimes U_{>\ell})\ket{n,q}=\ket{n}\otimes\ket{x_{\le\ell}(n_<,q_<),y(n_>,q_>)}_\sigma$.
No fermionic sign occurs, because both factors preserve occupation numbers and the identifications
are defined on occupation basis vectors.  $P^\pm_\ell$ selects $x_{b_\ell}=\pm1$, and on
$x_{b_\ell}=-1$ the operator $V_\ell$ flips every $y_{b_j}$ by \eqref{eq:szflip}.  Hence
$C_\ell(U_{\le\ell}\otimes U_{>\ell})\ket{n,q}=\ket{n}\otimes\ket{x(n,q)}_\sigma=U\ket{n,q}$.

(b) Let $\Psi':=(U_{\le\ell}\otimes U_{>\ell})\Phi$, $\rho':=\rho_A(\Psi')=U_{\le\ell}\,\rho_{A'}(\Phi)\,U_{\le\ell}^\dagger$,
and $\rho:=\rho_A(U\Phi)=\rho_A(C_\ell\Psi')$.  Then $S(\rho')=S_{A'}(\Phi)$.  Since
$(P^s_\ell\otimes\one)C_\ell=P^s_\ell\otimes V_s$ with $V_+=\one$, $V_-=V_\ell$, and the partial trace over
$\cH_B$ is invariant under conjugation by unitaries on $\cH_B$,
\[
  P^s_\ell\rho P^s_\ell=P^s_\ell\,\Tr_{\cH_B}\bigl[(\one\otimes V_s)\ket{\Psi'}\bra{\Psi'}(\one\otimes V_s^\dagger)\bigr]P^s_\ell
  =P^s_\ell\rho'P^s_\ell,\qquad s=\pm .
\]
So $\mathcal{P}(\rho)=\mathcal{P}(\rho')$ for the pinching of \cref{lem:pinching} with $P^\pm=P^\pm_\ell$.
Hence $S(\rho)\le S(\mathcal{P}(\rho))=S(\mathcal{P}(\rho'))\le S(\rho')+\ln2$, and symmetrically.

(c) Let $(-1)^{N_c}\Phi=\pi\Phi$ and $\prod_{k\le\ell}\taux{k}\Phi=Q\Phi$.  By \cref{thm:duality}(c) for the
chain $\{1,\dots,\ell\}$, $U_{\le\ell}^\dagger\sxb{\ell}U_{\le\ell}=\prod_{k\le\ell}\taux{k}(-1)^{n_k}$.  Under the
identification $(-1)^{N_c}=(-1)^{N_{\le\ell}}\otimes(-1)^{N_{>\ell}}$, and $U_{>\ell}$ commutes with
$(-1)^{N_{>\ell}}$.  Hence
\[
  (\sxb{\ell}\otimes\one)\Psi'=(U_{\le\ell}\otimes U_{>\ell})\bigl(Q(-1)^{N_{\le\ell}}\otimes\one\bigr)\Phi
  =Q\pi\,(\one\otimes(-1)^{N_{>\ell}})\Psi' .
\]
So $(P^\pm_\ell\otimes\one)\Psi'=(\one\otimes R^\pm)\Psi'$ with the complementary projections
$R^\pm:=\frac12(\one\pm Q\pi(-1)^{N_{>\ell}})$ on $\cH_B$, which commute with $V_\ell$.  Therefore
$U\Phi=C_\ell\Psi'=(\one\otimes(R^++V_\ell R^-))\Psi'$.  Here $R^++V_\ell R^-$ is a unitary on $\cH_B$,
so $\rho_A(U\Phi)=\rho'$ and $S_A(U\Phi)=S_{A'}(\Phi)$.  For $\Phi=\phi\otimes\ket{q}$, $\Phi$ is an eigenvector of
$\prod_{k\le\ell}\taux{k}$, and $\rho_{A'}(\Phi)=\rho^\phi_{\le\ell}\otimes\ket{q_{\le\ell}}\bra{q_{\le\ell}}$.

(d) Take $N=2$, $\ell=1$, $n=(0,0)$.  The vector $\Phi=\ket{0,0}_c\otimes2^{-1/2}(\ket{+1}+\ket{-1})_{\tau_1}\otimes\ket{+1}_{\tau_2}$
is a product over $A'|B'$, while by \eqref{eq:xnq}
$U\Phi=\ket{0,0}\otimes2^{-1/2}(\ket{+}_{b_1}\ket{+}_{b_2}+\ket{-}_{b_1}\ket{-}_{b_2})$ has $S_A=\ln2$.
Conversely, $\Psi=\ket{0,0}\otimes\ket{\uparrow}_{b_1}\otimes\ket{+}_{b_2}$ is a product over $A|B$, while
$U^\dagger\Psi=\ket{0,0}_c\otimes2^{-1/2}(\ket{+1,+1}+\ket{-1,-1})_{\tau}$ has $S_{A'}=\ln2$.
\end{proof}

\begin{remark}[Interior intervals]\label{rem:interior}
For an interval $A=[\ell_1+1,\ell_2]$ (sites and links) with $1\le\ell_1<\ell_2<N$, applying (a) at
$\ell_1$ and, inside the right factor (which is again a chain with boundary link $b_N$), at $\ell_2$ writes
$U=C_{\ell_1}(\one\otimes C'_{\ell_2})(U_{\le\ell_1}\otimes U_{(\ell_1,\ell_2]}\otimes U_{>\ell_2})$: a local unitary
followed by two controlled gates.  The gate $C'_{\ell_2}$ has its control $\sxb{\ell_2}$ in $A$ and its
target in the right complement, so by the argument of (b) it changes $S_A$ by at most $\ln2$.  The gate
$C_{\ell_1}$ has its control $\sxb{\ell_1}$ in the left complement.  Tracing out the complement, it changes the
reduced state of $A$ from $p_+\omega_++p_-\omega_-$ to $p_+\omega_++p_-V\omega_-V^\dagger$, where
$p_\pm\omega_\pm$ are the parts with $\sxb{\ell_1}=\pm1$ and $V=\prod_{\ell_1<k\le\ell_2}\szb{k}$.  Concavity of
$S$ and the bound $S(p_+\omega'_++p_-\omega'_-)\le p_+S(\omega'_+)+p_-S(\omega'_-)+\ln2$ for density matrices
$\omega'_\pm$ (proved as in \cref{lem:pinching}) show that this changes $S_A$ by at most $\ln2$.  Hence
$\abs{S_A(U\Phi)-S_{A'}(\Phi)}\le2\ln2$, with the Jordan--Wigner partial trace on both sides.  Numerically
the value $2\ln2$ is attained ($N=3$, $A$ consisting of site $2$ and link $b_2$).  We do not use this.
For interior intervals, odd elements of the algebra of $A$ carry the Jordan--Wigner string of the
left part.  For states of definite total fermion parity, however, the partial trace over the
complement is block diagonal with respect to the parity of $A$ and coincides with the fermionic
reduced state.
\end{remark}

\subsection{Sector averaging}\label{sec:model:sector}

For $q\in\{\pm1\}^N$ and $\psi\in\cF_N$ we write $\psi_q(t):=\Gamma(\e^{-\ii t\hq})\psi$.  $\E_q$
denotes expectation with respect to the uniform probability measure on $\{\pm1\}^N$, i.e.\ $q_1,\dots,q_N$
i.i.d.\ Bernoulli$(1/2)$ with values $\pm1$.  For a gauge-invariant $O$ the operator $U^\dagger OU$
commutes with every $\taux{j}$.  It therefore leaves every joint eigenspace $\cF^c_N\otimes\ket{q}$
invariant, and
\begin{equation}\label{eq:Oq}
  U^\dagger OU=\sum_qO_q\otimes\ket{q}\bra{q}
\end{equation}
for unique operators $O_q$ on $\cF^c_N$.  We set $O_q(t):=\Gamma(\e^{\ii t\hq})O_q\Gamma(\e^{-\ii t\hq})$.

\begin{corollary}[Exact sector averaging]\label{cor:sector}
Let $\psi\in\cF_N$ be a unit vector, $\Psi_0:=\ket{\uparrow\cdots\uparrow}_z\otimes\psi$ (we do not
distinguish $\LinkS_N\otimes\cF_N$ from $\cF_N\otimes\LinkS_N$), and $\Psi(t):=\e^{-\ii tH}\Psi_0$.
\begin{enumerate}[label=(\alph*)]
\item One has
\[
  U^\dagger\Psi_0=2^{-N/2}\sum_q\psi\otimes\ket{q},\qquad
  U^\dagger\Psi(t)=2^{-N/2}\sum_q\e^{\ii ht\sum_jq_j}\,\psi_q(t)\otimes\ket{q}.
\]
In particular $\norm{P_q\Psi_0}^2=2^{-N}$ for every $q$: the charges are exactly i.i.d.\
Bernoulli$(1/2)$, and every sector carries the same matter state $\psi$.
\item For every gauge-invariant $O$ and $t\in\R$,
\begin{equation}\label{eq:sectoravg}
  \langle\Psi(t),O\,\Psi(t)\rangle=\E_q\,\langle\psi_q(t),O_q\,\psi_q(t)\rangle .
\end{equation}
More generally, for gauge-invariant $O_1,\dots,O_m$ and $t_1,\dots,t_m\in\R$,
\[
  \langle\Psi_0,O_1(t_1)\cdots O_m(t_m)\Psi_0\rangle=\E_q\langle\psi,O_{1,q}(t_1)\cdots O_{m,q}(t_m)\psi\rangle,
\]
where $O(t)=\e^{\ii tH}O\e^{-\ii tH}$.  The constants $-h\sum_jq_j$ drop out.
\item \emph{Sector-diagonal observables.}  The following are gauge invariant, with the stated $O_q$:
\begin{itemize}
\item every function $F(n_1,\dots,n_N)$ of the occupations, with $O_q=F(n_1,\dots,n_N)$
independent of $q$; in particular $n_x$ and $\one[N_X=k]$.  Hence the distribution of $N_X$ in
$\Psi(t)$ is $\Prob_{\Psi(t)}(N_X=k)=\E_q\langle\psi_q(t),\one[N_X=k]\psi_q(t)\rangle$, the
$q$-average of the per-sector distributions;
\item every polynomial in the dressed fermions $\cdr_x,\cdr_x^\dagger$, with $O_q$ the same
polynomial in $c_x,c_x^\dagger$; in particular $f_x^\dagger(\prod_{x\le k<y}\szb{k})f_y\mapsto c_x^\dagger c_y$;
\item the electric field and energy: $\sxb{j}\mapsto(\prod_{k\le j}q_k)(-1)^{N_{\le j}}$ and
$\sxb{j-1}\sxb{j}\mapsto q_j(-1)^{n_j}$, together with their products.
\end{itemize}
For a Fock state $\psi=\ket{n}$ with occupied set $\{s_1<\dots<s_M\}$, write $u_q(t):=\e^{-\ii t\hq}$.
Then $\langle\psi_q(t),c_x^\dagger c_y\psi_q(t)\rangle=\sum_{a=1}^M\overline{u_q(t)_{xs_a}}\,u_q(t)_{ys_a}$, and
\begin{equation}\label{eq:stringdet}
  \langle\psi_q(t),(-1)^{N_{\le j}}\psi_q(t)\rangle=\det\bigl[(u_q(t)^*R_ju_q(t))_{s_as_b}\bigr]_{a,b=1}^M,
  \qquad R_j:=\diag\bigl((-1)^{\one[x\le j]}\bigr)_{x=1}^N .
\end{equation}
\item \emph{The gauge field $\szb{j}$ is not sector diagonal.}  It anticommutes with $G_j$ and
$G_{j+1}$ (with $G_N$ only, if $j=N$), and
\begin{equation}\label{eq:loschmidt}
  \langle\Psi(t),\szb{j}\Psi(t)\rangle=\E_q\Bigl[\e^{\ii ht(\sum_kq_k-\sum_kq'_k)}\,
  \langle\psi_{q'}(t),\psi_q(t)\rangle\Bigr],\qquad q':=q^{[j]},
\end{equation}
where $q^{[j]}$ is $q$ with $q_j$ and $q_{j+1}$ flipped ($q_N$ only, if $j=N$), so that
$\sum_kq_k-\sum_kq'_k=2(q_j+q_{j+1})$ (resp.\ $2q_N$).  For Fock $\psi$,
$\langle\psi_{q'}(t),\psi_q(t)\rangle=\det[(u_{q'}(t)^*u_q(t))_{s_as_b}]_{a,b=1}^M$.  Thus gauge-field
observables are Loschmidt echoes between sectors that differ by a local flip of the potential, and
the sector-dependent phase in \eqref{eq:loschmidt} cannot be omitted.
\end{enumerate}
\end{corollary}

\begin{proof}
(a) $\ket{\uparrow}=2^{-1/2}(\ket{+}+\ket{-})$, so $\ket{\uparrow\cdots\uparrow}_z=2^{-N/2}\sum_x\ket{x}_\sigma$.
For fixed $n$, $U^\dagger(\ket{n}\otimes\ket{x}_\sigma)=\ket{n,q(n,x)}$ with $q(n,x)$ given by
\eqref{eq:qofx}, and $x\mapsto q(n,x)$ is a bijection of $\{\pm1\}^N$.  Hence
$U^\dagger(\ket{n}\otimes\ket{\uparrow\cdots\uparrow}_z)=2^{-N/2}\sum_q\ket{n,q}$, and the claim
follows by linearity in $\psi$.  The time evolution follows from \cref{thm:duality}(e).
(b) By (a) and \eqref{eq:Oq},
\[
  \langle\Psi(t),O\Psi(t)\rangle=2^{-N}\sum_{q,q'}\e^{\ii ht(\sum_kq_k-\sum_kq'_k)}
  \bigl\langle\psi_{q'}(t)\otimes\ket{q'},(O_q\psi_q(t))\otimes\ket{q}\bigr\rangle,
\]
and only $q'=q$ contributes.  For the multi-time formula,
$U^\dagger O_i(t_i)U=\sum_qO_{i,q}(t_i)\otimes\ket{q}\bra{q}$ by \cref{thm:duality}(e); the phases cancel
because $O_i$ is block diagonal.
(c) Gauge invariance and the $O_q$ follow from \cref{thm:duality}(c),(d) and \cref{lem:gauss}.  The
determinant formulas follow from $\psi_q(t)=c^\dagger(u_qe_{s_1})\cdots c^\dagger(u_qe_{s_M})\ket0$
(with $u_q=u_q(t)$), from $(-1)^{N_{\le j}}=\Gamma(R_j)$, and from
\[
  \bigl\langle c^\dagger(v_1)\cdots c^\dagger(v_M)\ket0,\,c^\dagger(w_1)\cdots c^\dagger(w_M)\ket0\bigr\rangle
  =\det[\langle v_a,w_b\rangle]_{a,b=1}^M .
\]
(d) By \cref{thm:duality}(c), $U^\dagger\szb{j}U=\tauz{j}\tauz{j+1}$ maps $\cF^c_N\otimes\ket{q}$ onto
$\cF^c_N\otimes\ket{q^{[j]}}$.  Hence, by (a),
\[
  \langle\Psi(t),\szb{j}\Psi(t)\rangle
  =2^{-N}\sum_q\e^{-\ii ht\sum_kq^{[j]}_k}\,\e^{\ii ht\sum_kq_k}\,
  \bigl\langle\psi_{q^{[j]}}(t),\psi_q(t)\bigr\rangle .
\]
The determinant formula is the overlap formula used in (c).
\end{proof}

\begin{remark}[Superposition versus ensemble]\label{rem:superposition}
By \eqref{eq:sectoravg}, the dynamics of every gauge-invariant observable in the coherent
superposition $\Psi_0$ coincides with the dynamics in the \emph{ensemble} of definite-charge states
$\psi\otimes\ket{q}$ (in the dual picture), $q$ i.i.d.  The superposed state and the ensemble are
distinguished only by observables that are not gauge invariant, such as $\szb{j}$ in
\eqref{eq:loschmidt}, and by entanglement-type quantities (\cref{lem:cutlocal},
\cref{prop:entropies}).
\end{remark}

\begin{lemma}[Sector laws of product initial states]\label{lem:productlaw}
Let $\chi=\bigotimes_{j=1}^N\chi_j\in\LinkS_N$ be a product state of the links.  Write
$\chi(x):=\prod_j\langle x_{b_j}|\chi_j\rangle$ for its $\sigma^x$-amplitudes, and
$\Psi_0=\chi\otimes\psi$.
\begin{enumerate}[label=(\alph*)]
\item For every $\psi=\sum_n\psi(n)\ket{n}$, $U^\dagger\Psi_0=\sum_q\psi^{(q)}\otimes\ket{q}$ with
$\psi^{(q)}:=\sum_n\chi(x(n,q))\psi(n)\ket{n}$.
\item \emph{Fock matter.}  If $\psi=\ket{n}$, then $\psi^{(q)}=\chi(x(n,q))\ket{n}$.  The sector law
$\mu(q):=\norm{P_q\Psi_0}^2=\abs{\chi(x(n,q))}^2$ is the image of the product law
$\bigotimes_j\bigl(\abs{\langle+|\chi_j\rangle}^2\delta_{+1}+\abs{\langle-|\chi_j\rangle}^2\delta_{-1}\bigr)$
of $x$ under the bijection $x\mapsto q$, $q_j=(-1)^{n_j}x_{b_{j-1}}x_{b_j}$, $x_{b_0}=1$.  Every
sector in the support carries the matter state $\ket{n}$, up to a phase, and for every gauge-invariant
$O$, $\langle\Psi(t),O\Psi(t)\rangle=\sum_q\mu(q)\langle n,O_q(t)\,n\rangle$.
\item \emph{$\sigma^z$-product gauge states.}  If $\chi=\ket{s}_z$ with $D:=\{j:s_j=-1\}$, then for every
$\psi$
\[
  U^\dagger\Psi_0=2^{-N/2}\sum_q\varepsilon_D(q)\,(W_D\psi)\otimes\ket{q},\qquad
  \varepsilon_D(q):=\prod_{j\in D}\prod_{k\le j}q_k,\quad W_D:=\prod_{j\in D}(-1)^{N_{\le j}} .
\]
So the charges are i.i.d.\ Bernoulli$(1/2)$, and \cref{cor:sector}(b),(c) hold with $\psi$ replaced
by $W_D\psi$.  For Fock $\psi$, $W_D\psi=\pm\psi$, and the gauge-invariant dynamics is identical to
that of $\ket{\uparrow\cdots\uparrow}_z\otimes\psi$.
\item \emph{$\sigma^x$-product gauge states.}  If $\chi=\ket{\xi}_\sigma$, $\xi\in\{\pm1\}^N$, and
$\psi=\ket{n}$, then $U^\dagger\Psi_0=\ket{n}\otimes\ket{q(n,\xi)}$, a single (definite-charge) sector with
$q_j=(-1)^{n_j}\xi_{b_{j-1}}\xi_{b_j}$.
\item \emph{Mixed product states and the dimer law.}  Let $N=2M$ be even.  Let the odd links
$b_{2k-1}$ be in $\sigma^z$-eigenstates and the even links $b_{2k}$ in the $\sigma^x$-eigenstates
$\ket{\xi_{2k}}$, $\xi_{2k}\in\{\pm1\}$, and let $\psi=\ket{n}$.  Then $(q_{2k-1})_{k=1}^M$ are i.i.d.\ uniform
under $\mu$, and $q_{2k}=d_k\,q_{2k-1}$ with the deterministic signs
$d_k:=(-1)^{n_{2k-1}+n_{2k}}\xi_{2k-2}\xi_{2k}$, $\xi_0:=1$.  In particular, for
\[
  \chi_{2k-1}=\ket{\uparrow}_z,\qquad \chi_{2k}=\ket{(-1)^k}\ \ (\sigma^x\text{-eigenvalue }(-1)^k),\qquad
  n_{2k-1}=0,\quad n_{2k}=1,
\]
for $k=1,\dots,M$,
one has, with all amplitudes positive,
\begin{equation}\label{eq:dimerstate}
  U^\dagger\Psi_0=2^{-M/2}\sum_{\eta\in\{\pm1\}^M}\ket{n}\otimes\ket{\eta_1,\eta_1,\eta_2,\eta_2,\dots,\eta_M,\eta_M}.
\end{equation}
The sector law is the \emph{random dimer law}: i.i.d.\ equal pairs $(q_{2k-1},q_{2k})=(\eta_k,\eta_k)$,
$\eta_k$ i.i.d.\ uniform.  The same CDW matter state sits in every sector.
\end{enumerate}
\end{lemma}

\begin{proof}
(a) Expand $\chi=\sum_x\chi(x)\ket{x}_\sigma$ and use $U^\dagger(\ket{n}\otimes\ket{x}_\sigma)=\ket{n,q(n,x)}$
together with the bijection $x\leftrightarrow q$ for fixed $n$.
(b) This is (a) for $\psi=\ket{n}$.  The probability of $x$ under the product law is
$\abs{\chi(x)}^2$, and the formula for gauge-invariant $O$ follows as in \cref{cor:sector}(b).
(c) $\ket{\downarrow}=\sigma^x\ket{\uparrow}$, so $\Psi_0=\bigl(\prod_{j\in D}\sxb{j}\bigr)(\ket{\uparrow\cdots\uparrow}_z\otimes\psi)$.
By \cref{thm:duality}(c),
$U^\dagger\prod_{j\in D}\sxb{j}U=\prod_{j\in D}\bigl(\prod_{k\le j}\taux{k}\bigr)\,\prod_{j\in D}(-1)^{N_{\le j}}$;
apply it to \cref{cor:sector}(a).  The sign $\varepsilon_D(q)$ has modulus one and $W_D$ is a $q$-independent
unitary.  It acts on Fock states as a sign.
(d) This is (a) with $\chi(x)=\delta_{x,\xi}$.
(e) On the support of $\chi(x)$, $x_{b_{2k}}=\xi_{2k}$, while $y_k:=x_{b_{2k-1}}$ are free, with
$\abs{\chi(x)}^2=2^{-M}$.  By \eqref{eq:qofx}, $q_{2k-1}=(-1)^{n_{2k-1}}\xi_{2k-2}\,y_k$ and
$q_{2k}=(-1)^{n_{2k}}y_k\,\xi_{2k}$.  So $q_{2k-1}$ is a fixed sign times $y_k$, i.i.d.\ uniform, and
$q_{2k}q_{2k-1}=d_k$.  For the specific state, $\langle y|\uparrow\rangle=2^{-1/2}>0$ for both values of
$y$, and $\xi_{2k-2}\xi_{2k}=(-1)^{k-1}(-1)^k=-1$ (also for $k=1$, since $\xi_0=1$, $\xi_2=-1$).
Hence $d_k=(-1)^{0+1}\cdot(-1)=1$, $q_{2k-1}=q_{2k}=(-1)^{k-1}y_k=:\eta_k$, and $y\mapsto\eta$ is a
bijection, which gives \eqref{eq:dimerstate}.
\end{proof}

\begin{remark}\label{rem:productlaws}
\leavevmode
\begin{enumerate}[label=(\alph*)]
\item \cref{lem:productlaw}(c) is the precise form of the observation in the Supplementary Material of
\cite{SKKM2017} that every $\sigma^z$-product spin state is gauge equivalent to the polarized one:
the sector law stays i.i.d.\ uniform, and the matter state changes by the $q$-independent unitary
$W_D$.
\item For the biased $\sigma^x$-product states $\chi_j=\cos\theta\ket{+}+\sin\theta\ket{-}$ with Fock
matter, (b) shows that $q_j=(-1)^{n_j}x_{b_{j-1}}x_{b_j}$ with $x$ i.i.d.\ and
$\Prob(x_{b_j}=1)=\cos^2\theta$.  This is a block factor of an i.i.d.\ sequence: $q_j$ is a function of
$(x_{b_{j-1}},x_{b_j})$, so $q_j$ and $q_k$ are independent for $\abs{j-k}\ge2$.  For $N\ge2$ the charges are
independent only if $\cos^2\theta\in\{0,\tfrac12,1\}$: with $m:=2\cos^2\theta-1$ one has
$\E[(-1)^{n_1}q_1]=m$, $\E[(-1)^{n_2}q_2]=m^2$ and $\E[(-1)^{n_1+n_2}q_1q_2]=\E[x_{b_2}]=m$, and $m\neq m^3$
unless $m\in\{0,\pm1\}$.  For $\cos^2\theta=\tfrac12$ the law is the i.i.d.\ law of \cref{cor:sector}; for
$\cos^2\theta\in\{0,1\}$ it is a point mass (a single sector, \cref{lem:productlaw}(d)).
\item The localization statements of \cref{sec:main} are proved for the i.i.d.\ law.  An upper bound of
the form $\E_q[F(q)]\le B$ with $F\ge0$ transfers to every sector law $\mu$ with $\mu(q)\le K\,2^{-N}$ for all
$q$, as $\sum_q\mu(q)F(q)\le KB$; e.g.\ $K=2$ for the open chain (\cref{prop:variants}(i)).  Lower bounds,
such as the positivity of the infinite-time imbalance in \cref{thm:imbalance}, do not transfer in this
way.  The dimer law \eqref{eq:dimerstate} has density $2^{N/2}$ on its support, which is unbounded in $N$;
in infinite volume the two laws are mutually singular.
\Cref{sec:dimer} shows that for $0<2\abs{h}<J$ the dimer law leads to transport that is at least diffusive.
Disorder-free localization is a property of the sector law induced by the initial state, not of the
Hamiltonian alone.
\end{enumerate}
\end{remark}

\begin{remark}[The clean sector is ballistic]\label{rem:gaussproj}
Let $\Psi_0\in P_q\cH_N$ for one fixed $q$.  This is the case, for instance, for the strictly
gauge-invariant subspace $G_j\equiv1$ of the usual lattice-gauge-theory formulation (no background
charges), or for the definite-charge states $\ket{\xi}_\sigma\otimes\ket{n}$ of
\cref{lem:productlaw}(d).  Then, by \cref{thm:duality}(e), the matter evolves under the single,
deterministic one-body operator $\hq$, and nothing is averaged.  Write $\one_N:=(1,\dots,1)$.  Then $h_{\one_N}=-J\Delta_N+2h\one$ is the clean hopping
operator shifted by a constant, and the localization properties of \cref{sec:main} fail in this sector:
\begin{enumerate}[label=(\roman*)]
\item \emph{No decay of the eigenfunction correlator.}  The eigenvectors of $h_{\one_N}$ are
\[
  \varphi_k(x)=\sqrt{\tfrac{2}{N+1}}\,\sin\tfrac{\pi kx}{N+1},\qquad k=1,\dots,N,
\]
and $\sin\frac{\pi kN}{N+1}=\pm\sin\frac{\pi k}{N+1}$.  Since
$\sum_{k=1}^N\sin^2\frac{\pi k}{N+1}=\frac{N+1}{2}$,
\[
  \Qc^{[1,N]}_{\one_N}(1,N)=\sum_{k=1}^N\abs{\varphi_k(1)\varphi_k(N)}=1\qquad\text{for every }N .
\]
This is the largest possible value, since always $\Qc\le1$.  So the deterministic analogue of
Hypothesis~\ref{hyp:EC} fails in this sector.
\item \emph{No memory of a density wave.}  Let $N$ be even, $\Gamma:=\diag((-1)^x)$, and
$a_k:=\langle\varphi_k,\Gamma\varphi_k\rangle$.  Since $\Gamma\Delta_N\Gamma=-\Delta_N$, we have
$h_{\one_N}\Gamma+\Gamma h_{\one_N}=4h\Gamma$.  Taking the expectation in $\varphi_k$ gives
$(E_k-2h)a_k=0$, where $E_k=2h-2J\cos(\pi k/(N+1))\neq2h$ because $N+1$ is odd.  Hence $a_k=0$ for
all $k$.  Take the CDW Fock state $\psi$ ($n_x=1$ iff $x$ is even) in this sector, e.g.\
$\Psi_0=\ket{\xi}_\sigma\otimes\psi$ with $\xi_{b_j}=\prod_{k\le j}(-1)^{n_k}$.  Its imbalance
$I(t):=\frac2N\sum_x(-1)^x\langle\Psi(t),n_x\Psi(t)\rangle$ equals
$\frac1N\Tr(\Gamma\e^{-\ii th_{\one_N}}\Gamma\e^{\ii th_{\one_N}})$, because $\Tr\Gamma=0$.  Its infinite-time
average $\lim_{T\to\infty}\frac1T\int_0^TI(t)\,\dd t$ is $\frac1N\sum_ka_k^2=0$, since the spectrum is simple.  Compare \cref{thm:imbalance}, where
the $q$-average makes the time average strictly positive for every $h\neq0$.
\item \emph{Ballistic transport.}  On $\ell^2(\Z)$, let $h_\Z:=-J\Delta+2h$ and let $X$ be the
position operator.  Since $h_\Z$ is bounded and of finite range, $\e^{-\ii th_\Z}$ maps finitely
supported vectors into the domain of every polynomial in $X$.  On such vectors
$\e^{\ii th_\Z}X\e^{-\ii th_\Z}=X+tV$, with $V=\ii[h_\Z,X]=-\ii J(S_+-S_-)$, $(S_\pm u)(x)=u(x\pm1)$,
because $V$ commutes with $h_\Z$.  Hence
$\langle\delta_y,\e^{\ii th_\Z}(X-y)^2\e^{-\ii th_\Z}\delta_y\rangle=t^2\norm{V\delta_y}^2=2J^2t^2$
for all $t$.  Consequently there is no bound
$\sup_t\abs{\langle\delta_x,\e^{-\ii th_\Z}\delta_y\rangle}\le F(\abs{x-y})$ with
$\sum_dd^2F(d)^2<\infty$.  By \eqref{eq:heisenberg}, the dressed-fermion propagator restricted to
$P_{\one_N}\cH_N$ is $\e^{-\ii th_{\one_N}}$, and at fixed $t$ its entries at distance $\gg\abs t(2J+2\abs h)$
from the ends converge to those of $\e^{-\ii th_\Z}$ as $N\to\infty$ (finite propagation speed,
\cref{lem:propagation}(b)).  This is why the zero-velocity Lieb--Robinson bound of \cref{prop:LR} can only
hold on average over sectors, and not in operator norm uniformly in $N$.
\end{enumerate}
For the initial states of \cref{cor:sector}, whose gauge part $\ket{\uparrow\cdots\uparrow}_z$ is translation
invariant, disorder-free localization is therefore produced by the uniform average over the $2^N$
sectors, almost all of which carry a disordered potential.  It fails in the clean sector $q=\one_N$; in
particular it is invisible in the strictly gauge-invariant subspace, and no bound that is uniform over all
sectors (such as an operator-norm Lieb--Robinson bound) can hold.  Single sectors with a disordered charge
configuration do localize.  For instance, for the definite-charge states $\ket{\xi}_\sigma\otimes\ket{n}$
with $\xi$ uniformly distributed, $\Sfull$ obeys an area law in expectation (\cref{thm:arealaw}(c)).  In
that case, however, the disorder is already present in the initial gauge configuration $\xi$.
\end{remark}

\begin{remark}[Numerical verification]\label{rem:EDcheck}
The script
\begin{center}
\texttt{code/model\_duality\_check.py}\quad(output: \texttt{code/model\_duality\_check.out})
\end{center}
builds $H$, the $G_j$, the $f_j$ (Jordan--Wigner, ordering \eqref{eq:occbasis}) and the link Pauli
matrices on $\cH_N$ from scratch.  It builds $U$ as the explicit matrix
$(\one\otimes\mathrm{Had}^{\otimes N})\,\Pi$, where $\Pi$ is the permutation matrix of
$(n,q)\mapsto(n,x(n,q))$ and $\mathrm{Had}$ is the Hadamard matrix.  It checks, for $N=2,\dots,6$ and
$J=1$, $h=0.7$:
\begin{itemize}
\item every identity of \cref{thm:duality}(a)--(d), and \eqref{eq:gausslaw1d}, to $\le10^{-14}$;
\item the spectral identity of \cref{thm:duality}(b) (full spectrum of dimension $4^N\le4096$ against
the union of the free-fermion sector spectra), max $\abs{\Delta E}\le1.0\times10^{-13}$;
\item the open and ring identities of \cref{prop:variants}, as operator identities and as spectra.
Sector by sector for the ring ($N=3,4$), twist $w$ has error $\le6\times10^{-15}$, while the
alternative Jordan--Wigner-type twist $w(-1)^{N_f+1}$ has error $2.0$;
\item the cut decomposition of \cref{lem:cutlocal}(a) for all $\ell$, to $3\times10^{-17}$.  It also
checks \cref{lem:cutlocal}(b),(c) on random vectors ($N=3,4$; equality in (c) to $2\times10^{-15}$)
and the sharpness examples of (d).  Along the dynamics from $\ket{\uparrow\cdots\uparrow}_z\otimes$CDW
($N=6$, $\ell=3$, $h=J$) the defect $\abs{S_A(\Psi(t))-S_{A'}(U^\dagger\Psi(t))}$ reaches $0.496$ (at $t=2$),
so it is a genuine effect, of the order of the bound $\ln2\approx0.693$;
\item the dynamics of \cref{cor:sector}: full exact diagonalization of the LGT against
\eqref{eq:sectoravg}, \eqref{eq:stringdet}, \eqref{eq:loschmidt} for $N=4,5,6$, $h\in\{0.4,1.3\}$,
$t\le10$, for the CDW and for random non-Fock $\psi$.  The observables are $n_x$, all dressed
bilinears, $\sxb{j}$, $\szb{j}$ and the distribution of $N_{\le\lfloor N/2\rfloor}$; the maximal
deviation is $\le10^{-14}$.  Dropping the phase in \eqref{eq:loschmidt} gives an error of $0.78$
($N=4$, $h=1.3$);
\item \cref{lem:productlaw}(b)--(e), including the dimer law \eqref{eq:dimerstate} for $N=2,4,6$, to
$10^{-15}$;
\item \cref{rem:gaussproj}: $a_k=0$ to $3\times10^{-14}$ for even $N\le200$,
$\Qc_{\one_N}^{[1,N]}(1,N)=1$ to $10^{-14}$ for $N\le300$, and the ballistic law on a chain of
length $1201$ (relative error $3\times10^{-15}$ for $t\le100$).
\end{itemize}
\end{remark}
